\subsection{由局部到整体的过渡}

命题 8. 设 A 是一个环，m 是 A 的一个极大理想，M 是一个 A-模。若存在 A 的理想 a，使得 m 是 A 的包含 a 的唯一极大理想，且 aM = 0，则典范同态 \(M \rightarrow M_{m}\) 是双射。

此时 A/a 是以 m/a 为极大理想的局部环；可以把 M 视为一个 (A/a)-模；对一切 \(s \in A - m\)，s 在 A/a 中的典范像是可逆的，因而 M 的位似 \(x \mapsto sx\) 是双射，这源于把 \(M_{m}\) 定义为泛问题的解（§2 第 2 节）；由此即得命题。

特别地，若存在 \(k \geqslant 0\) 使 \(\mathfrak{m}^k\mathbf{M} = 0\)，则同态 \(\mathbf{M} \to \mathbf{M}_m\) 是双射（§1 第 1 节命题 1 的推论）。

命题 9. 设 A 是一个环，m 是 A 的一个极大理想，M 是一个 A-模，k 是一个整数 \(\geqslant0\)。典范同态 \(M\to M_{m}/m^{k}M_{m}\) 是满射，其核为 \(m^{k}M\)，并且给出 \(M/m^{k}M\) 到 \(M_{m}/m^{k}M_{m}\) 上的一个同构。

由于情形 \(k=0\) 是平凡的，设 \(k\geqslant1\)。由命题 8，典范同态 \(\mathbf{M}/\mathfrak{m}^{k}\mathbf{M}\to(\mathbf{M}/\mathfrak{m}^{k}\mathbf{M})_{\mathfrak{m}}\) 是双射。另一方面，\((\mathbf{M}/\mathfrak{m}^{k}\mathbf{M})_{\mathfrak{m}}\) 典范地等同于 \(\mathbf{M}_{\mathfrak{m}}/(\mathfrak{m}^{k}\mathbf{M})_{\mathfrak{m}}\)（§2 第 4 节定理 1），因而 \((\mathfrak{m}^{k}\mathbf{M})_{\mathfrak{m}}=\mathfrak{m}^{k}\mathbf{M}_{\mathfrak{m}}\)（§2 第 7 节命题 18 的推论），从而存在 \(\mathbf{M}/\mathfrak{m}^{k}\mathbf{M}\) 到 \(\mathbf{M}_{\mathfrak{m}}/\mathfrak{m}^{k}\mathbf{M}_{\mathfrak{m}}\) 上的一个同构，它把元素 \(x\in\mathbf{M}\) 的类映为 \(x/1\) 的类。

推论. 设 A 是一个环，\(m_{1}, m_{2}, \ldots, m_{n}\) 是 A 的互不相同的极大理想，M 是一个 A-模，\(k_{1}, k_{2}, \ldots, k_{n}\) 是整数 \(\geqslant 0\)。从 M 到 \(\prod_{i=1}^{n} \mathbf{M}_{\mathfrak{m}_{i}}/\mathfrak{m}_{i}^{k_{i}} \mathbf{M}_{\mathfrak{m}_{i}}\) 的典范同态是满射，其核为 \(\left(\bigcap_{i=1}^{n} \mathfrak{m}_{i}^{k_{i}}\right) \mathbf{M}\)。

这容易由命题 9 以及 §1 第 2 节命题 6 得出，因为诸 \(\mathfrak{m}_{i}^{k_{i}}\) 两两互素（§1 第 2 节命题 3）。

在本节余下部分，A 表示一个环，\(\Omega(A)\)（或 \(\Omega\)）表示 \(A\) 的极大理想的集合。

命题 10. A-模 \(\bigoplus_{m\in\Omega}A_m\)，即诸 \(A_m\)（其中 \(m\in\Omega\)）的直和，是忠实平坦的。

每个 \(A_{m}\) 都是平坦 A-模（\(\S\) 2 第 4 节定理 1），因而 \(E = \bigoplus_{m \in \Omega} A_{m}\) 是平坦的（第 I 章 \(\S\) 2 第 3 节命题 2）。此外，对 A 的每个极大理想 m，\(mA_{m}\) 是 \(A_{m}\) 的唯一极大理想，故 \(mA_{m} \neq A_{m}\)，由此推出 \(mE \neq E\)，从而 E 是忠实平坦的（第 I 章 \(\S\) 3 第 1 节命题 1 (d)）。

定理 1. 设 M、N 是两个 A-模，\(u: \mathbf{M} \to \mathbf{N}\) 是一个 A-同态，且对每个 \(m \in \Omega\)，设 \(u_m: \mathbf{M}_m \to \mathbf{N}_m\) 是相应的 \(\mathbf{A}_m\)-同态（§2 第 2 节注 5）。为使 \(u\) 是单射（相应地，满射、双射、零射），必须且只需对每个 \(m \in \Omega\)，\(u_m\) 是单射（相应地，满射、双射、零射）。

说对每个 \(m \in \Omega\) \(u_m\) 是单射（相应地，满射、双射、零射），等价于说同态 \(\bigoplus_{m} u_m: \bigoplus_{m} M_m \to \bigoplus_{m} N_m\) 具有同样的性质。但 \(\bigoplus_{m} M_m = M \otimes_A E\)，\(\bigoplus_{m} N_m = N \otimes_A E\)，且 \(\bigoplus_{m} u_m = u \otimes 1\)，其中 \(E = \bigoplus_{m} A_m\)；由于 \(E\) 是忠实平坦的（命题 10），定理由第 I 章 §3 第 1 节命题 1 (c) 及命题 2 得出。

推论 1. 设 M 是一个 A-模，N 是 M 的一个子模，x 是 M 的一个元素。为使 \(x \in N\)，必须且只需对每个 \(m \in \Omega\)，x 在 \(M_{m}\) 中的典范像属于 \(N_{m}\)。

设 \(\bar{x}\) 是 \(x\) 在 M/N 中的类；说 \(x \in N\)，意指从 A 到 M/N 的 A-线性映射 \(u: \alpha \mapsto \alpha \bar{x}\) 是零射。现在，\((\mathrm{M} / \mathrm{N})_{\mathfrak{m}}\) 等同于 \(\mathrm{M}_{\mathfrak{m}} / \mathrm{N}_{\mathfrak{m}}\)（§2 第 4 节定理 1），而 \(u_{\mathfrak{m}}: \mathrm{A}_{\mathfrak{m}} \to \mathrm{M}_{\mathfrak{m}} / \mathrm{N}_{\mathfrak{m}}\) 等同于映射 \(\lambda \mapsto \lambda \bar{x}_{\mathfrak{m}}\)，其中 \(\bar{x}_{\mathfrak{m}}\) 是模 \(\mathrm{N}_{\mathfrak{m}}\) 意义下的类，即 \(x\) 在 \(\mathrm{M}_{\mathfrak{m}}\) 中的典范像的类。由于由定理 1，关系 \(u = 0\) 等价于对一切 m 有 \(u_{\mathfrak{m}} = 0\)，这就证明了推论。

推论 2. 设 \(\mathbf{M}\) 是一个 A-模，且对每个 \(\mathfrak{m} \in \Omega\)，设 \(f_{\mathfrak{m}}\) 是典范映射 \(\mathbf{M} \to \mathbf{M}_{\mathfrak{m}}\)。同态 \(x \mapsto (f_{\mathfrak{m}}(x))\)（从 \(\mathbf{M}\) 到 \(\prod_{\mathfrak{m} \in \Omega} \mathbf{M}_{\mathfrak{m}}\)）是单射。

在 N = 0 的情形应用推论 1 可见，关系 x = 0 等价于 \(f_{m}(x) = 0\) 对一切 \(m \in \Omega\) 成立。

推论 3. (i) 设 \(\mathfrak{b}\) 是 \(\mathbf{A}\) 的一个理想，a 是 \(\mathbf{A}\) 的一个元素。为使 \(a \in \mathfrak{b}\)，必须且只需对每个 \(\mathfrak{m} \in \Omega\)，\(a\) 在 \(\mathbf{A}_{\mathfrak{m}}\) 中的典范像属于 \(\mathfrak{b}\mathbf{A}_{\mathfrak{m}}\)。

\begin{enumerate}
\def\labelenumi{(\roman{enumi})}
\setcounter{enumi}{1}
\tightlist
\item
  特别地，设 \(b\) 与 \(c\) 是 \(A\) 的两个元素。为使 \(c\) 是 \(b\) 的倍元，必须且只需对每个 \(m \in \Omega\)，\(c\) 在 \(A_m\) 中的典范像是 \(b\) 的典范像的倍元。
\end{enumerate}

由于 \(\mathfrak{bA}_{m} = \mathfrak{b}_{m}\)（§2 第 7 节命题 18 的推论），(i) 是推论 1 的特殊情形；(ii) 由把 (i) 应用于理想 Ab 而得。

推论 4. 设 A 是一个整环，K 是其分式域，M 是一个无挠 A-模，使得 M 典范地等同于 \(\mathbf{K} \otimes_{\mathbf{A}} \mathbf{M}\) 的一个 A-子模。那么，对每个 \(\mathfrak{m} \in \Omega\)，\(\mathbf{M}_{\mathfrak{m}}\) 典范地等同于 \(\mathbf{K} \otimes_{\mathbf{A}} \mathbf{M}\) 的一个 A-子模，且 \(\mathbf{M} = \bigcap_{\mathfrak{m} \in \Omega} \mathbf{M}_{\mathfrak{m}}\)。

由于 M 等同于 \(K \otimes_{A} M\) 的一个子模，故 \(M_{m}\) 等同于 \(A_{m}\)-模 \((\mathbf{K} \otimes_{\mathbf{A}} \mathbf{M})_{\mathfrak{m}} = \mathbf{K}_{\mathfrak{m}} \otimes_{\mathbf{A}} \mathbf{M}\) 的一个子模（§2 第 4 节定理 1）；由于 \(K_{m} = K\)，我们立即看出 \(M_{m}\) 是无挠的；此外，图的交换性

\[
\begin{array}{c} \mathbf {M} \longrightarrow \mathbf {K} \otimes_ {\mathbf {A}} \mathbf {M} \\ \Big \downarrow \qquad \qquad \qquad \qquad \qquad \qquad \qquad \qquad \Big \downarrow \\ \mathbf {M} _ {\mathfrak {m}} \longrightarrow (\mathbf {K} \otimes_ {\mathbf {A}} \mathbf {M}) _ {\mathfrak {m}} \end{array}
\]

证明典范映射 \(M \rightarrow M_{m}\) 是单射。于是该推论由把推论 1 应用于 A-模 \(K \otimes_{A} M\) 及其子模 M 而得。

特别地，对每个整环 A，

\[
\mathbf {A} = \bigcap_ {\mathfrak {m} \in \Omega} \mathbf {A} _ {\mathfrak {m}}.\tag{2}
\]

推论 5. 设 A 是一个环。A-模 \(\mathbf{A}^n\) 的每个由 \(n\) 个元素组成的生成系都是 \(\mathbf{A}^n\) 的一个基。

设 \((e_i)_{1\leqslant i\leqslant n}\) 是 \(\mathbf{A}^n\) 的典范基，\((x_i)_{1\leqslant i\leqslant n}\) 是 \(\mathbf{A}^n\) 的一个由 \(n\) 个元素组成的生成系，\(u\colon \mathbf{A}^n\to \mathbf{A}^n\) 是满足 \(u(e_i) = x_i\)（对 \(1\leqslant i\leqslant n\)）的 A-线性映射。由假设 \(u\) 是满射，需要证明 \(u\) 是单射。凭定理 1，这可以立即化为 \(\mathbf{A}\) 是局部环的情形；若 \(\mathfrak{m}\) 是 \(\mathbf{A}\) 的极大理想，则 \(1\otimes x_i\)（\(1\leqslant i\leqslant n\)）在 \((\mathrm{A / m})^n\) 中构成自由 \((\mathrm{A / m})\)-模 \((\mathrm{A / m})^n\) 的一个生成系；由于 \(\mathrm{A / m}\) 是域，该生成系是 \((\mathrm{A / m})^n\) 的一个基；由于 \(\mathbf{A}^n\) 是自由 A-模，我们由命题 5 推出 \((x_i)\) 是 \(\mathbf{A}^n\) 的一个基。

命题 11. 设 M 是一个 A-模，N 是一个有限生成 A-模，\(u: \mathbf{M} \to \mathbf{N}\) 是一个同态。为使 \(u\) 是满射，必须且只需对每个 \(\mathfrak{m} \in \Omega\)，同态 \(\mathbf{M} / \mathfrak{m}\mathbf{M} \to \mathbf{N} / \mathfrak{m}\mathbf{N}\)（由 \(u\) 取商而得）是满射。

由定理 1，为使 \(u\) 是满射，必须且只需 \(u_{\mathfrak{m}}: \mathrm{M}_{\mathfrak{m}} \to \mathrm{N}_{\mathfrak{m}}\) 对一切 \(\mathfrak{m} \in \Omega\) 是满射。由于 \(\mathrm{A}_{\mathfrak{m}}\) 是局部环且 \(\mathrm{N}_{\mathfrak{m}}\) 是有限生成 \(\mathrm{A}_{\mathfrak{m}}\)-模，这等价于说由取商得到的同态 \(u_{\mathfrak{m}}': \mathrm{M}_{\mathfrak{m}} / \mathrm{mM}_{\mathfrak{m}} \to \mathrm{N}_{\mathfrak{m}} / \mathrm{mN}_{\mathfrak{m}}\) 是满射（第 2 节命题 4 的推论 1）；而 \(\mathrm{M}_{\mathfrak{m}} / \mathrm{mM}_{\mathfrak{m}}\)（相应地 \(\mathrm{N}_{\mathfrak{m}} / \mathrm{mN}_{\mathfrak{m}}\)）等同于 \(\mathrm{M} / \mathrm{mM}\)（相应地 \(\mathrm{N} / \mathrm{mN}\)）（命题 9），由此即得命题。

命题 12. 设 E、F、G 是三个 A-模，\(v: \mathbf{G} \to \mathbf{F}\) 与 \(u: \mathbf{E} \to \mathbf{F}\) 是同态。假设 E 是有限表示的。为使存在同态 \(w: \mathbf{E} \to \mathbf{G}\) 使 \(u\) 分解为 \(\mathbf{E} \xrightarrow{w} \mathbf{G} \xrightarrow{v} \mathbf{F}\)，必须且只需对每个 \(m \in \Omega\)，存在同态 \(w^m: \mathbf{E}_m \to \mathbf{G}_m\) 使 \(u_m: \mathbf{E}_m \to \mathbf{F}_m\) 分解为 \(\mathbf{E}_m \xrightarrow{w^m} \mathbf{G}_m \xrightarrow{v_m} \mathbf{F}_m\)。

存在满足上述命题的 w 等价于下面的性质：u 属于映射

\[
r = \operatorname{Hom} (1 _ {\mathbf {E}}, v): \operatorname{Hom} _ {\mathbf {A}} (\mathbf {E}, \mathbf {G}) \rightarrow \operatorname{Hom} _ {\mathbf {A}} (\mathbf {E}, \mathbf {F}).
\]

现在，\((\mathrm{Hom}_{\mathbf{A}}(\mathrm{E},\mathbf{F}))_{\mathfrak{m}}\)（相应地 \((\mathrm{Hom}_{\mathbf{A}}(\mathrm{E},\mathbf{G}))_{\mathfrak{m}}\)）典范地等同于 \(\mathrm{Hom}_{\mathbf{A}_{\mathfrak{m}}}(\mathrm{E}_{\mathfrak{m}},\mathrm{F}_{\mathfrak{m}})\)（相应地 \(\mathrm{Hom}_{\mathbf{A}_{\mathfrak{m}}}(\mathrm{E}_{\mathfrak{m}},\mathrm{G}_{\mathfrak{m}})\)）（§2 第 7 节命题 19 (i)），u 在 \((\mathrm{Hom}_{\mathbf{A}}(\mathrm{E},\mathbf{F}))_{\mathfrak{m}}\) 中的典范像等同于 \(u_{m}\)，\(r_{m}\) 等同于 \(\mathrm{Hom}_{\mathbf{A}_{\mathfrak{m}}}(1_{\mathbf{E}_{\mathfrak{m}}},v_{\mathfrak{m}})\)，而 \(P_{m}\) 等同于 \(r_{m}\) 的像。于是该命题由把定理 1 的推论 1 应用于 \(\mathrm{Hom}_{\mathbf{A}}(\mathrm{E},\mathbf{F})\) 及其子模 P 而得。

推论 1. 设 M 是一个 A-模，N 是 M 的一个子模，使得 M/N 是有限表示的。为使 N 是 M 的一个直因子，必须且只需对每个 \(m \in \Omega\)，\(N_{m}\) 是 \(M_{m}\) 的一个直因子。

说 N 是 M 的直因子，意指 M/N 上的恒同同态分解为 \(M/N \xrightarrow{w} M \xrightarrow{\phi} M/N\)，其中 \(\phi\) 是典范同态，w 是一个同态（《代数》第 II 章 §1 第 9 节命题 14）；由于 \((\mathrm{M}/\mathrm{N})_{\mathrm{m}} = \mathrm{M}_{\mathrm{m}}/\mathrm{N}_{\mathrm{m}}\) 且 \(\phi_{m}\) 是典范同态 \(M_{m} \to M_{m}/N_{m}\)，故由命题 12 容易得出推论。

推论 2. 设 M 是一个有限生成的自由 A-模，N 是 M 的一个子模且是有限生成的自由 A-模。为使 N 是 M 的一个直因子，必须且只需对所有 m ∈ Ω，mN = N ∩ (mM)。

按定义 M/N 是有限表示的；另一方面，\(N_{m}\) 与 \(M_{m}\) 是有限生成的自由 \(A_{m}\)-模。为使 \(N_{m}\) 是 \(M_{m}\) 的一个直因子，必须且只需典范映射 \(N_{m}/mN_{m} \rightarrow M_{m}/mM_{m}\) 是单射（第 2 节命题 6）；这等价于说典范映射 \(N/mN \rightarrow M/mM\) 是单射（命题 9），而由于其核为 \((N \cap mM)/mN\)，这就证明了推论。

命题 12（相应地，其推论 1）尤其将应用于 A 是 Noether 环且 E（相应地，M/N）是有限生成 A-模的情形（第 I 章 §2 第 8 节引理 8）。

